\section*{Chapter \ref{ch:s2:harvesting-the-variance-risk-premium} --- Harvesting the Variance Risk Premium}

\begin{solution}{exo:s2:harvesting-the-variance-risk-premium:1}
$0.8 \times 0.20 \times \sqrt{21/252} = 0.8 \times 0.20 \times 0.289 = 4.6\%$ of the index.
\end{solution}

\begin{solution}{exo:s2:harvesting-the-variance-risk-premium:2}
$0.16^2 \times 1.3 = 0.0333$, an implied volatility of $\sqrt{0.0333} = 18.2\%$ against an expected 16\%.
\end{solution}

\begin{solution}{exo:s2:harvesting-the-variance-risk-premium:3}
The put is out of the money: it expires worthless in most months (92\%) and pays the full premium, but in a crash its loss is many times the premium, a
skewness of $-7.2$; the variance swap's P\&L is more even across months, so its mean is a larger share of its volatility.
\end{solution}

\begin{solution}{exo:s2:harvesting-the-variance-risk-premium:4}
A monthly standard deviation is $10\%/\sqrt{12} = 2.9\%$; 7.7 of them are 22.2\% of capital.
\end{solution}

\begin{solution}{exo:s2:harvesting-the-variance-risk-premium:5}
The planted crashes start from calm markets, when implied variance and therefore the strike are low, so the book is largest exactly when the crash
begins. Inverse-strike sizing would help if large losses followed high volatility, as when a volatile period continues rather than a calm one
breaking.
\end{solution}

\begin{solution}{exo:s2:harvesting-the-variance-risk-premium:6}
Investors pay for protection against variance, which rises when they are poorest. The average of VIX minus realised volatility includes the months when
realised far exceeded the VIX; the premium is measured on variance, where those months weigh more, and over long samples that include crashes.
\end{solution}

\begin{solution}{exo:s2:harvesting-the-variance-risk-premium:7}
At 30\% the worst month costs 66.5\% and the capital survives; at 50\% it costs 110.8\% and the capital is gone. The worst month is 7.68 monthly standard
deviations, so the whole capital goes at a target of $1/(7.68/\sqrt{12}) = 45\%$.
\end{solution}

\begin{solution}{exo:s2:harvesting-the-variance-risk-premium:8}
Five calm years say nothing about the crash that the strategy is paid for bearing; the synthetic book's history before its worst month also had a
Sharpe ratio above 1.5. Size the fund by a stress (a month like February 2020 or October 2008 applied to today's positions), not by its history.
\end{solution}

\begin{solution}{pb:s2:harvesting-the-variance-risk-premium:1}
\begin{enumerate}
\item Selling options or variance to earn the premium; selling options against a position held, most often calls.
\item The variance risk premium as realised variance minus a variance swap rate synthesised from options, on five indexes and 35 stocks.
\item VIX 19.5 against realised 15.4 on average, higher in 84\% of months, by 4.1 points; variance 31\% higher.
\item February 2020, September and August 2008, March 2025, July 2011, June 2002.
\item Stochastic variance with leverage, jumps and three crashes; implied variance 1.3 times expected variance, a falling smile.
\item 0.69, 1.24, 0.40, 0.08 and 0.09; skewness from $-2.3$ to $-12.5$.
\item The premium is paid for bearing crashes: small gains most months, large losses rarely.
\item PutWrite 7.2\% a year (Sharpe 0.57), BuyWrite 6.0\% (0.48); $-27.6\%$ and $-28.5\%$ in autumn 2008, $-19.8\%$ and $-21.3\%$ in early 2020.
\item The gain of a delta-hedged option relative to its premium, which measures the volatility premium.
\item Delta-hedged long index options underperform zero, less away from the money, more when volatility is high.
\item 0.53\% of the index a month, about an eighth of the premium; the hedge removes the index's direction.
\item Fewer hedge trades and more path dependence: lower costs, noisier P\&L.
\item \textbf{Named result.} The short-variance book's Sharpe ratio before its worst month is 1.57 at any size; the worst month costs 11.1\%, 22.2\%,
44.3\% and 88.6\% of capital at 5\%, 10\%, 20\% and 40\% volatility, and more than the capital at 60\%.
\item The crash is a single rare event; five calm years contain no information about its size.
\item By a stress on today's positions, so that the worst plausible month is survivable.
\item It sized up in calm markets, where the crashes began.
\item Crash years in the sample, executable option prices, margin calls, and a direct comparison with the underlying.
\item Put writing wins most often (92\% of months); the hedged straddle has the best Sharpe ratio.
\item Volatility targeting cuts risk after volatility rises; here the losses come as volatility rises, faster than a target can react.
\item Insurance against variance.
\end{enumerate}
\end{solution}

\begin{solution}{iq:s2:harvesting-the-variance-risk-premium:1}
The difference between the price of variance (implied) and the variance then realised; it is positive on average because variance is high in bad times
and investors pay to hedge it.
\end{solution}

\begin{solution}{iq:s2:harvesting-the-variance-risk-premium:2}
The straddle's gamma loss is about $\tfrac12\Gamma(\Delta S)^2$, large for a 5\% gap, and implied volatility rises too (a vega loss); re-hedge the new
delta at the open, check the stress limit, and cut if the book is beyond it rather than wait for a rebound.
\end{solution}

\begin{solution}{iq:s2:harvesting-the-variance-risk-premium:3}
With stress limits on the loss in historical and hypothetical crashes applied to current positions, vega and gamma limits by tenor, a cap on the
premium collected relative to capital, and liquidity limits on how fast the book can be cut.
\end{solution}

\begin{solution}{iq:s2:harvesting-the-variance-risk-premium:4}
Variance swaps take the premium across all strikes and have a clean, convex payoff with crash exposure in the tail; straddles take it at the money
with delta hedging and path dependence; out-of-the-money puts take the richest part of the skew with the most concentrated crash risk.
\end{solution}

\begin{solution}{iq:s2:harvesting-the-variance-risk-premium:5}
Positions and prices, a surface for each underlying, Greeks by option, aggregation by underlying, a hedging rule (time or bands), order generation in
the hedge instrument, and reconciliation of hedge fills against the model.
\end{solution}

\begin{solution}{iq:s2:harvesting-the-variance-risk-premium:6}
The hedged position's change is $\tfrac12\Gamma(dS)^2 + \Theta\,dt$; with $\Theta = -\tfrac12\Gamma S^2\sigma_i^2$ (Black--Scholes at implied volatility) and
$(dS)^2 = S^2\sigma_r^2\,dt$, the P\&L is $\tfrac12\Gamma S^2(\sigma_r^2 - \sigma_i^2)\,dt$: a long option earns when realised exceeds implied, a short one
when implied exceeds realised, weighted by gamma.
\end{solution}
